% Texto restaurado a partir do PDF entregue.
\section{Motores de Passo NEMA 23, Passo de 0,9$^\circ$ e Sele\c{c}\~ao da Corrente $I_{RMS}$}
\label{sec:nema23_irms}

\subsection{NEMA 23: geometria e implica\c{c}\~oes mec\^anicas}

NEMA 23 define as dimens\~oes da \emph{flange} do motor
(aproximadamente \SI{57}{\milli\meter} $\times$
\SI{57}{\milli\meter}) e o padr\~ao de fura\c{c}\~ao de montagem, mas n\~ao
define torque, corrente nominal ou passo el\'etrico. Assim, motores com
essa carca\c{c}a podem variar em comprimento do corpo, in\'ercia do rotor,
resist\^encia e indut\^ancia por fase, al\'em do torque de reten\c{c}\~ao.
No projeto deste trabalho, os tr\^es eixos utilizam motores NEMA 23 por
combinarem: (i) facilidade de fixa\c{c}\~ao em estruturas CNC;
(ii) bom compromisso entre torque e tamanho; e
(iii) ampla disponibilidade de \emph{drivers} compat\'iveis, como o
TMC5160.

\subsection{Passo el\'etrico: 1,8$^\circ$ vs 0,9$^\circ$}

O \^angulo de passo el\'etrico corresponde \`a rota\c{c}\~ao do eixo a cada
passo inteiro sem \emph{microstepping}. Os valores mais usuais s\~ao
$1{,}8^\circ$ e $0{,}9^\circ$, o que leva a
\begin{equation}
    N_{\text{passos/rev}} = \frac{360^\circ}{\theta_{\text{passo}}}.
\end{equation}
Assim, motores de $1{,}8^\circ$ possuem 200 passos por volta, enquanto os
de $0{,}9^\circ$ possuem 400 passos por volta.

Na pr\'atica, o passo de $0{,}9^\circ$ oferece maior resolu\c{c}\~ao
mec\^anica por volta e reduz o \emph{ripple} posicional em baixas
velocidades. Em contrapartida, exige taxa de pulsos mais alta para a
mesma velocidade angular. Para um fator de microstepping $M$, vale
\begin{equation}
    \text{RPS} = \frac{f_{\text{STEP}}}{N \cdot M}
    \quad \Rightarrow \quad
    f_{\text{STEP}} = \text{RPS} \cdot N \cdot M.
\end{equation}
Portanto, a mesma rota\c{c}\~ao por segundo em um motor de $0{,}9^\circ$
demanda o dobro da frequ\^encia de pulsos em rela\c{c}\~ao a um motor de
$1{,}8^\circ$.

\subsection{Microstepping, suavidade e taxa de passos}

O \emph{microstepping} interpola correntes aproximadamente senoidais nas
fases do motor, reduzindo ru\'ido, vibra\c{c}\~ao e rugosidade de
movimento. Entretanto, a resolu\c{c}\~ao efetiva de comando continua
limitada pela frequ\^encia de pulsos \texttt{STEP} dispon\'ivel no
controlador. Al\'em disso, o torque incremental associado a cada
micropasso \'e menor do que o torque de passo inteiro, raz\~ao pela qual a
opera\c{c}\~ao com cargas din\^amicas exige correntes compat\'iveis com a
acelera\c{c}\~ao desejada e com a curva torque versus velocidade do motor.

\subsection{Par\^ametros el\'etricos e modelo de primeira ordem}

Motores de passo bipolares s\~ao especificados por corrente por fase
$I_{\text{rated}}$, resist\^encia $R$ e indut\^ancia $L$. Em primeira
aproxima\c{c}\~ao, a din\^amica de corrente da fase pode ser escrita como
\begin{equation}
    \frac{di}{dt} \approx \frac{V_{\text{bus}} - v_{\text{chopper}}}{L},
    \qquad
    \tau = \frac{L}{R}.
\end{equation}
No TMC5160, o controle por corrente aplica \emph{chopping} para manter o
valor de $I_{RMS}$ desejado, usando tens\~ao de barramento elevada para
acelerar a resposta de corrente contra a indut\^ancia do enrolamento. O
aquecimento cresce aproximadamente com $P \propto I_{RMS}^2 R$, de modo
que trabalhar com fra\c{c}\~oes moderadas da corrente nominal reduz a
dissipa\c{c}\~ao sem necessariamente comprometer o torque requerido.

\subsection{Torques relevantes}

\begin{itemize}
\item \textbf{Torque de reten\c{c}\~ao} ($T_{\text{hold}}$): torque
  est\'atico m\'aximo com corrente nominal e eixo travado.
\item \textbf{Torque de detente} ($T_{\text{det}}$): ondula\c{c}\~ao
  mec\^anica intr\'inseca do rotor sem corrente aplicada.
\item \textbf{Torque din\^amico}: decresce com a velocidade el\'etrica
  devido \`a limita\c{c}\~ao de $di/dt$ e \`a for\c{c}a
  contraeletromotriz; por isso, barramentos mais altos ajudam a preservar
  torque em rota\c{c}\~oes elevadas.
\end{itemize}

\subsection{Resumo para os motores usados}

\begin{itemize}
\item \textbf{JK57HM76-2804}: NEMA 23, corrente nominal de
  aproximadamente \SI{2.8}{\ampere} por fase e torque de reten\c{c}\~ao em
  torno de \SI{1.8}{\newton\meter}. \'E adequado aos eixos que exigem
  maior for\c{c}a est\'atica.
\item \textbf{23HM8430 0,9$^\circ$}: NEMA 23, corrente nominal de
  aproximadamente \SI{3.0}{\ampere} e torque de reten\c{c}\~ao pr\'oximo de
  \SI{1.5}{\newton\meter}. O passo fino nativo favorece resolu\c{c}\~ao e
  suavidade em baixas velocidades.
\end{itemize}

\subsection{Fundamenta\c{c}\~ao te\'orica para a sele\c{c}\~ao de $I_{RMS}$}

Como primeira aproxima\c{c}\~ao, o torque de reten\c{c}\~ao pode ser usado
para estimar a constante de torque do motor:
\begin{equation}
    k_T \approx \frac{T_{\text{hold}}}{I_{\text{rated}}}
    \qquad [\si{\newton\meter\per\ampere}].
\end{equation}
Com isso, uma estimativa da corrente m\'inima pr\'atica para superar o
torque de detente pode ser escrita como
\begin{equation}
    I_{\min} \approx \alpha \frac{T_{\text{det}}}{k_T},
\end{equation}
em que $\alpha$ \'e um fator de seguran\c{c}a escolhido entre 1,5 e 2,5
para compensar atrito e irregularidades din\^amicas.

Assumindo $T_{\text{det}} \approx \SI{0.06}{\newton\meter}$ como valor
t\'ipico para motores NEMA 23, obt\'em-se a Tabela~\ref{tab:irms-estimado}.

\begin{table}[H]
    \centering
    \caption{Estimativa da corrente m\'inima pr\'atica ($I_{RMS}$).}
    \label{tab:irms-estimado}
    \setlength{\tabcolsep}{4pt}\footnotesize
    \begin{tabularx}{\textwidth}{lccXX}
        \toprule
        Motor & $I_{\text{rated}}$ [A] & $T_{\text{hold}}$ [N$\cdot$m] & $k_T$ [N$\cdot$m/A] & $I_{\min}$ [A] \\
        \midrule
        JK57HM76-2804 & 2,8 & 1,8 & 0,643 & 0,14--0,23 \\
        23HM8430 0,9$^\circ$ & 3,0 & 1,5 & 0,50 & 0,18--0,30 \\
        \bottomrule
    \end{tabularx}
\end{table}

Do ponto de vista t\'ermico e de robustez, adotou-se no firmware um
perfil de corrente reduzida para os testes de bancada, com
$I_{\text{RUN}} \approx \SI{0.28}{\ampere}$ e
$I_{\text{HOLD}} \approx \SI{0.12}{\ampere}$ no motor JK57HM76-2804, e
$I_{\text{RUN}} \approx \SI{0.30}{\ampere}$ com
$I_{\text{HOLD}} \approx \SIrange{0.12}{0.15}{\ampere}$ no 23HM8430.
Esses valores representam apenas uma fra\c{c}\~ao da corrente nominal,
reduzindo expressivamente a pot\^encia dissipada sem preju\'izo
significativo sobre o torque necess\'ario aos movimentos previstos neste
prot\'otipo CNC. Caso o processo venha a exigir maior acelera\c{c}\~ao ou
maior esfor\c{c}o mec\^anico, o ajuste de $I_{\text{RUN}}$ pode ser feito
dinamicamente no TMC5160 por meio do registrador
\texttt{IHOLD\_IRUN}.
